% Publication-style paper (PDF). Build: latexmk -pdf main.tex
% Figures: ../output/*.png relative to latex/

\documentclass[11pt,a4paper]{article}

\usepackage[T1]{fontenc}
\usepackage[utf8]{inputenc}
\usepackage{lmodern}
\usepackage{microtype}
\usepackage{geometry}
\geometry{margin=1in}

\usepackage{amsmath,amssymb}
\usepackage{booktabs}
\usepackage{graphicx}
\usepackage{float}
\usepackage{caption}
\usepackage{subcaption}
\usepackage{xcolor}

% Links / refs: subtle light-green PDF borders (Overleaf-style), not big framed callouts
\usepackage[
  colorlinks=false,
  pdfborder={0 0 0.35},
  linkbordercolor={0.62 0.80 0.70},
  citebordercolor={0.62 0.80 0.70},
  urlbordercolor={0.55 0.76 0.65},
  filebordercolor={0.62 0.80 0.70},
  bookmarks=true,
  bookmarksopen=true,
]{hyperref}
\usepackage{xurl}

\usepackage[
  backend=biber,
  style=authoryear,
  maxcitenames=2,
  uniquename=init,
]{biblatex}
\addbibresource{references.bib}

\newcommand{\fig}[2]{%
  \begin{center}
    \IfFileExists{#1}{%
      \includegraphics[width=0.92\linewidth]{#1}%
    }{%
      \begin{minipage}{0.9\linewidth}
        \small\sffamily
        \fbox{\begin{minipage}{0.98\linewidth}
          \path{#1}\\[0.35em] #2
        \end{minipage}}
      \end{minipage}%
    }%
  \end{center}%
}

\title{\textbf{Predicting Conflict in Africa:\\[0.3em]
\large A Machine Learning Approach Using Environmental Stressors}}
\author{%
  Andrés Felipe Camacho\thanks{Corresponding: \texttt{afcamachob@uchicago.edu}} \and
  Pablo Hernández Pedraza \and
  Agustín Eyzaguirre\\[0.8em]
  \normalsize Masters in Computational Analysis and Public Policy (MSCAPP)\\
  University of Chicago
}
\date{\today}

\begin{document}
\maketitle

\begin{abstract}
\noindent
This paper investigates the predictability of violent conflict in Africa using machine learning on satellite-derived environmental and socioeconomic data for $50\mathrm{km}\times 50\mathrm{km}$ grid cells (2019--2023). We compare logistic regression, $k$-nearest neighbors, random forest, and neural networks, trained on climate, wealth, and historical conflict features to predict binary conflict occurrence. A conservative weighted ensemble achieves approximately $72\%$ recall and $52\%$ precision on the minority (conflict) class, illustrating both the value and the precision--recall tradeoffs of early-warning style models for policy and humanitarian audiences.
\end{abstract}

\section{Introduction}

Conflict remains a major barrier to development and stability, particularly where environmental stress and resource scarcity interact with inequality. Predicting where violence is more likely can help allocate aid and diplomatic attention---but models must balance catching true conflicts (recall) against false alarms (precision). Africa offers a demanding testbed: heterogeneous geography, uneven data quality, and strong class imbalance (most grid-years are non-conflict).

\paragraph{Research question.}
\textit{Can the occurrence of violence in a given geography be predicted using local socioeconomic and environmental data?}

\paragraph{Related work.}
Our project sits in a broader agenda of using machine learning for conflict prediction. \cite{hegre2024views} describe the \textbf{VIEWS} system (Uppsala/PRIO), which produces grid-level fatality forecasts using gradient-boosted trees and related methods. \cite{perry2013machines} discusses whether machine learning can support conflict prediction in policy contexts. For $k$-NN background, see \textcite{knnlecture}.

\section{Data and outcomes}

\subsection{Sources}
We combine: \textbf{UCDP} (georeferenced events); \textbf{ERA5} (climate); \textbf{Meta RWI} (wealth); \textbf{Hansen} forest change; \textbf{VIIRS} nighttime lights. Grid cells are $50\mathrm{km}$; labels indicate at least one conflict event in a grid-year.

\subsection{Descriptive patterns}
Figure~\ref{fig:ts} shows the time series of violent events in Africa (1989--2023); Figure~\ref{fig:cty} shows concentration across countries.

\begin{figure}[H]
  \centering
  \begin{subfigure}[t]{0.48\linewidth}
    \fig{../output/conflicts_africa_89-23.png}{Export from the EDA pipeline.}
    \caption{Conflicts over time}
    \label{fig:ts}
  \end{subfigure}\hfill
  \begin{subfigure}[t]{0.48\linewidth}
    \fig{../output/conflicts_africa_country.png}{}
    \caption{Top countries}
    \label{fig:cty}
  \end{subfigure}
  \caption{Descriptive conflict patterns (Africa)}
\end{figure}

\subsection{Features and sample}
The training file has $34{,}015$ grid-year rows (2019--2023). Features include climate (temperature, precipitation, pressure), forest loss, lights, RWI, lags, and accumulated past conflict. The outcome is severely imbalanced ($\sim 92\%$ non-conflict); we use SMOTE during training for several models.

\begin{figure}[H]
  \centering
  \fig{../output/Data_Africa.png}{}
  \caption{Spatial view of conflict aggregation and temperature anomalies (example year).}
  \label{fig:map}
\end{figure}

\section{Models}

We implemented four models on the same dataset and tuned hyperparameters via cross-validation to make comparisons fair. We focus on accuracy and recall (especially for the minority conflict class).

\subsection{Logistic regression (Elastic Net)}
Our baseline is logistic regression with elastic-net regularization:
\[
P(y_i = +1 \mid x_i) = \sigma(w^\top x_i + b) = \frac{1}{1 + e^{-(w^\top x_i + b)}}.
\]
We minimize the negative log-likelihood plus ElasticNet penalty
\[
\ell(w,b) = \sum_{i=1}^n \bigl[ -\log \sigma\bigl(y_i(w^\top x_i + b)\bigr)\bigr]
+ \lambda \Bigl[ \alpha \sum_{j=1}^p |w_j| + (1-\alpha)\sum_{j=1}^p w_j^2 \Bigr].
\]
Features are standardized so coefficients are comparable. We tune $\alpha \in \{0,0.25,0.5,0.75,1\}$ and $C \in \{0.001,0.01,0.1,1,10\}$ (with $\lambda = 1/C$). We selected $\alpha=0$ (ridge-only) and $C=10$ (weak penalty). SMOTE addresses class imbalance; without it, the model would be biased toward the majority class.

\begin{table}[H]
  \centering
  \small
  \caption{Logistic regression---test classification report (excerpt).}
  \begin{tabular}{lcccc}
    \toprule
    Class & Precision & Recall & F1 & Support \\
    \midrule
    0 & 0.98 & 0.86 & 0.91 & 6251 \\
    1 & 0.32 & 0.75 & 0.45 & 552 \\
    \midrule
    \multicolumn{4}{l}{Accuracy} & 0.85 \\
    \bottomrule
  \end{tabular}
\end{table}

\subsection{\texorpdfstring{$k$-nearest neighbors}{k-nearest neighbors}}
We assume neighbors in feature space behave similarly. We sort training points by distance to $x_i$, form the set $S_k$ of $k$ nearest neighbors, and predict by distance-weighted majority vote
\[
\hat{y}_i = \arg\max_c \sum_{j \in S_k} w_j \cdot \mathbb{I}_{y_j=c}, \qquad
w_j = \frac{1}{d_j} \quad \text{(distance weights)}.
\]
We compare Euclidean and cosine distance, odd $k \in \{3,5,7,9,11\}$, and uniform vs.\ distance weights. The best specification uses \textbf{cosine similarity}, \textbf{$k=3$}, and distance-weighted voting, with test accuracy $0.89$ and weighted recall $0.89$.

\begin{table}[H]
  \centering
  \small
  \caption{$k$-NN---test classification report (excerpt).}
  \begin{tabular}{lcccc}
    \toprule
    Class & Precision & Recall & F1 & Support \\
    \midrule
    0 & 0.97 & 0.90 & 0.93 & 6236 \\
    1 & 0.38 & 0.68 & 0.49 & 567 \\
    \midrule
    \multicolumn{4}{l}{Accuracy} & 0.88 \\
    \bottomrule
  \end{tabular}
\end{table}

\subsection{Random forest}
Random Forest averages $B$ trees trained on bootstrap samples with random feature subsets at each split; we use majority vote for binary labels and average tree votes for probability of conflict. We tuned \texttt{n\_estimators}, \texttt{max\_depth}, \texttt{min\_samples\_split}, \texttt{min\_samples\_leaf}, and \texttt{class\_weight} via randomized search with 5-fold stratified CV. The selected model uses $500$ trees, unlimited depth, and no class weighting (SMOTE already handles imbalance). Performance is strong on both classes; feature importance ranks \texttt{accumulated\_conflicts} and temperature-related variables highly.

\begin{table}[H]
  \centering
  \small
  \caption{Random forest---test classification report (excerpt).}
  \begin{tabular}{lcccc}
    \toprule
    Class & Precision & Recall & F1 & Support \\
    \midrule
    0 & 0.97 & 0.95 & 0.96 & 6251 \\
    1 & 0.55 & 0.70 & 0.62 & 552 \\
    \midrule
    \multicolumn{4}{l}{Accuracy} & 0.89 \\
    \bottomrule
  \end{tabular}
\end{table}

\subsection{Neural network}
We use a feedforward network with two hidden layers (64 and 32 units, ReLU) and sigmoid output
\[
h_1 = \mathrm{ReLU}(W_1 x + b_1),\quad
h_2 = \mathrm{ReLU}(W_2 h_1 + b_2),\quad
P(\mathrm{conflict}) = \sigma(W_3 h_2 + b_3) .
\]
Training uses binary cross-entropy, SGD with momentum $0.9$, learning rate $0.01$, batch size $128$, and 100 epochs. The model achieves high recall on the minority class at the cost of lower precision; this profile can suit early-warning settings where missing conflicts is costly.

\begin{table}[H]
  \centering
  \small
  \caption{Neural network---test classification report (excerpt).}
  \begin{tabular}{lcccc}
    \toprule
    Class & Precision & Recall & F1 & Support \\
    \midrule
    0 & 0.98 & 0.88 & 0.93 & 6251 \\
    1 & 0.38 & 0.80 & 0.51 & 552 \\
    \midrule
    \multicolumn{4}{l}{Accuracy} & 0.88 \\
    \bottomrule
  \end{tabular}
\end{table}

\subsection{Ensemble}
We combine predicted probabilities from all four models with nonnegative weights summing to one and threshold $0.5$. We evaluated several schemes (equal weights, RF-heavy, recall-focused, conservative, etc.). The \textbf{Conservative} ensemble (LR 30\%, KNN 25\%, RF 35\%, NN 10\%) balances accuracy and recall for deployment-oriented use: it achieves about $92.3\%$ weighted accuracy and $72.3\%$ recall on the conflict class with $52.0\%$ precision on that class, alongside strong ROC AUC relative to single models.

\begin{table}[H]
  \centering
  \small
  \caption{Ensemble schemes (test set; see notebook for full grid).}
  \begin{tabular}{lccccc}
    \toprule
    Scheme & Accuracy & Precision & Recall & F1 & ROC AUC \\
    \midrule
    Conservative & 0.923 & 0.520 & 0.723 & 0.605 & 0.929 \\
    RF-dominant & 0.927 & 0.540 & 0.705 & 0.611 & 0.937 \\
    Balanced & 0.919 & 0.501 & 0.734 & 0.595 & 0.933 \\
    Equal weights & 0.917 & 0.493 & 0.728 & 0.588 & 0.929 \\
    \bottomrule
  \end{tabular}
\end{table}

\section{Results}

\subsection{Confusion matrices}
Figures~\ref{fig:cm_lr}--\ref{fig:cm_nn} summarize predicted versus actual labels for each model. Off-diagonal mass shows where models confuse conflict and non-conflict cells; comparing panels highlights how tree-based and ensemble methods trade off false positives and false negatives relative to linear and local methods.

\begin{figure}[H]
  \centering
  \begin{subfigure}[t]{0.48\linewidth}
    \fig{../output/LR_confussion.png}{}
    \caption{Logistic regression}
    \label{fig:cm_lr}
  \end{subfigure}\hfill
  \begin{subfigure}[t]{0.48\linewidth}
    \fig{../output/KNN_confussion.png}{}
    \caption{\texorpdfstring{$k$-NN}{k-NN}}
    \label{fig:cm_knn}
  \end{subfigure}
\end{figure}

\begin{figure}[H]
  \centering
  \begin{subfigure}[t]{0.48\linewidth}
    \fig{../output/confusion_RF.png}{}
    \caption{Random forest}
    \label{fig:cm_rf}
  \end{subfigure}\hfill
  \begin{subfigure}[t]{0.48\linewidth}
    \fig{../output/confusion_NN.png}{}
    \caption{Neural network}
    \label{fig:cm_nn}
  \end{subfigure}
\end{figure}

\begin{figure}[H]
  \centering
  \fig{../output/RF_Importance.png}{}
  \caption{Random forest variable importance (impurity-based).}
\end{figure}

\subsection{ROC and precision--recall}
\begin{figure}[H]
  \centering
  \begin{subfigure}[t]{0.48\linewidth}
    \fig{../output/ROC_Models.png}{}
    \caption{ROC curves}
  \end{subfigure}\hfill
  \begin{subfigure}[t]{0.48\linewidth}
    \fig{../output/Precision_Recall.png}{}
    \caption{Precision--recall}
  \end{subfigure}
\end{figure}

\begin{table}[H]
  \centering
  \small
  \caption{Model comparison (weighted metrics; see notebook).}
  \begin{tabular}{lccccc}
    \toprule
    Metric & LR & KNN & RF & NN & Ensemble \\
    \midrule
    Class 0 recall & 0.861 & 0.910 & 0.950 & 0.895 & 0.941 \\
    Class 1 recall & 0.754 & 0.679 & 0.697 & 0.772 & 0.723 \\
    Weighted recall & 0.852 & 0.891 & 0.929 & 0.885 & 0.923 \\
    Weighted accuracy & 0.852 & 0.891 & 0.929 & 0.885 & 0.923 \\
    Weighted precision & 0.923 & 0.924 & 0.938 & 0.931 & 0.938 \\
    \bottomrule
  \end{tabular}
\end{table}

\section{Discussion}

Grid-level conflict prediction combines genuine signal with hard tradeoffs: high recall often implies many false positives, which can erode trust among practitioners. Ensembles help but do not remove the need for domain judgment and ethical safeguards if scores inform allocation or enforcement.

\section{Conclusion}

Environmental and socioeconomic features contain information about conflict risk at subnational scale, but models should be validated out-of-sample, communicated with uncertainty, and paired with governance discussion---especially where scores could be misused. Future work includes richer institutional covariates, better missing-data strategies for remote sensing, and tighter links to operational early-warning workflows.

\paragraph{Replication.}
Code and data pipeline:\linebreak
{\small\url{https://github.com/anfelipecb/conflict-prediction-ml}}. The project site hosts an interactive Kepler viewer under \texttt{docs/}.

\printbibliography

\end{document}
